\documentclass[12pt]{article}
\usepackage[utf8]{inputenc}
\usepackage[T1]{fontenc}
\usepackage[margin=1in]{geometry}
\usepackage{setspace}
\usepackage{xcolor}
\usepackage[colorlinks=true, urlcolor=blue!60!black, linkcolor=blue!60!black]{hyperref}
\setlength{\parindent}{2em}
\setlength{\parskip}{1.25em}
\renewcommand{\baselinestretch}{1.0}
\begin{document}
\begin{center}
{\LARGE \textit{Tennessee v. Lane}, 541 U.S. 509 (2004)\par}
\vspace{0.8em}
{\large Case Brief\par}
\end{center}
\vspace{1em}

\section{Memory Jogger}
Congress validly exercised its enforcement power under Section~5 of the Fourteenth Amendment to abrogate state sovereign immunity under Title~II of the Americans with Disabilities Act insofar as Title~II applies to the class of cases implicating the fundamental right of access to the courts (\href{https://www.law.cornell.edu/supremecourt/text/541/509}{Cornell LII}; \href{https://www.courtlistener.com/opinion/134743/tennessee-v-lane/}{CourtListener}).

\section{Detailed Case Facts}
In August 1998, George Lane and Beverly Jones---both paraplegics who used wheelchairs for mobility---filed a federal lawsuit against the State of Tennessee and several Tennessee counties, alleging past and ongoing violations of Title~II of the Americans with Disabilities Act of 1990 (ADA), 42 U.S.C. \S\S~12131--12165 (\href{https://www.law.cornell.edu/supremecourt/text/541/509}{541 U.S. at 513}). Title~II mandates that ``no qualified individual with a disability shall, by reason of such disability, be excluded from participation in or be denied the benefits of the services, programs or activities of a public entity, or be subjected to discrimination by any such entity,'' 42 U.S.C. \S~12132, and incorporates by reference the private damages remedies of Section~505 of the Rehabilitation Act of 1973, 29 U.S.C. \S~794a, via 42 U.S.C. \S~12133 (\href{https://www.law.cornell.edu/supremecourt/text/541/509}{541 U.S. at 513, 517}). Congress also enacted an explicit abrogation provision in 42 U.S.C. \S~12202 stating that a State ``shall not be immune under the eleventh amendment'' from a federal or state court action for an ADA violation (\href{https://www.law.cornell.edu/supremecourt/text/541/509}{541 U.S. at 518}).

Lane had been summoned to appear in a county courthouse in Polk County, Tennessee, to answer criminal charges arising from driving with a revoked license and possessing prescription medication (\href{https://commons.und.edu/cgi/viewcontent.cgi?article=1006&context=ndlr}{North Dakota Law Review}). The courtroom was located on the second floor of a building with no elevator (\href{https://www.law.cornell.edu/supremecourt/text/541/509}{541 U.S. at 513--514}). For his first court appearance, Lane physically crawled up two flights of stairs to reach the courtroom (\href{https://www.law.cornell.edu/supremecourt/text/541/509}{541 U.S. at 514}). When he returned for a subsequent hearing, he refused to crawl up the stairs a second time or to be carried up the stairs by court officers; he was arrested and jailed for failure to appear (\href{https://www.law.cornell.edu/supremecourt/text/541/509}{541 U.S. at 514}). As Chief Justice Rehnquist noted in dissent, the county court held a preliminary hearing in the first-floor library and later offered to relocate subsequent proceedings to an accessible courthouse in a nearby town (\href{https://www.law.cornell.edu/supremecourt/text/541/509}{541 U.S. at 543 n.4}).

Beverly Jones, a certified court reporter, alleged that architectural barriers in multiple Tennessee county courthouses prevented her from gaining access to courtrooms, causing her to lose both employment and the opportunity to participate in the judicial process (\href{https://www.law.cornell.edu/supremecourt/text/541/509}{541 U.S. at 514}). Both plaintiffs sought monetary damages and equitable relief against the State and the defendant counties (\href{https://www.law.cornell.edu/supremecourt/text/541/509}{541 U.S. at 514}).

\section{Procedural History}
The State of Tennessee moved to dismiss the action in federal district court, asserting sovereign immunity under the Eleventh Amendment (\href{https://www.law.cornell.edu/supremecourt/text/541/509}{541 U.S. at 514}). The District Court denied the motion without opinion, and Tennessee took an interlocutory appeal under the collateral order doctrine (\href{https://www.law.cornell.edu/supremecourt/text/541/509}{541 U.S. at 514 \& n.1}). The United States intervened on appeal to defend the constitutionality of Title~II's abrogation of state sovereign immunity (\href{https://www.law.cornell.edu/supremecourt/text/541/509}{541 U.S. at 514}).

On April 28, 2000, the United States Court of Appeals for the Sixth Circuit held the appeal in abeyance pending the Supreme Court's decision in \textit{Board of Trustees of the University of Alabama v. Garrett}, 531 U.S. 356 (2001), which subsequently held that the Eleventh Amendment bars private money-damages suits against States under Title~I of the ADA (employment) while expressly reserving the question for Title~II (\href{https://www.law.cornell.edu/supremecourt/text/541/509}{541 U.S. at 514}). Following \textit{Garrett}, the en banc Sixth Circuit decided \textit{Popovich v. Cuyahoga County Court}, 276 F.3d 808 (6th Cir. 2002), holding that \textit{Garrett} bars Title~II damages claims premised on equal protection principles but permits Title~II damages claims grounded in due process principles (\href{https://www.law.cornell.edu/supremecourt/text/541/509}{541 U.S. at 514--515}).

Relying on \textit{Popovich}, a Sixth Circuit panel affirmed the denial of Tennessee's motion to dismiss, 40 Fed. Appx. 911 (6th Cir. 2002), explaining that respondents' claims rested on due process principles (\href{https://www.law.cornell.edu/supremecourt/text/541/509}{541 U.S. at 515}). After Tennessee petitioned for rehearing, the panel issued an amended opinion, 315 F.3d 680 (6th Cir. 2003), concluding that the Due Process Clause protects the right of access to the courts and that Congress's legislative record documented how physical barriers in courthouses denied disabled persons fundamental rights; the panel remanded for further factual development (\href{https://www.law.cornell.edu/supremecourt/text/541/509}{541 U.S. at 515--516}). The Supreme Court granted certiorari, 539 U.S. 941 (2003), heard oral argument on January 13, 2004, and affirmed on May 17, 2004 (\href{https://www.law.cornell.edu/supremecourt/text/541/509}{541 U.S. at 509, 516}; \href{https://www.oyez.org/cases/2003/02-1667}{Oyez}).

\section{Judicial Votes}
\textbf{Vote Split:} 5--4, affirming the Sixth Circuit (\href{https://www.law.cornell.edu/supremecourt/text/541/509}{541 U.S. at 512}; \href{https://supreme.justia.com/cases/federal/us/541/509/}{Justia}).

\begin{itemize}
    \item \textbf{Majority (5):} Justice John Paul Stevens delivered the opinion of the Court, joined by Justices Sandra Day O'Connor, David H. Souter, Ruth Bader Ginsburg, and Stephen G. Breyer.
    \item \textbf{Concurrences:} Justice David H. Souter filed a concurring opinion, joined by Justice Ruth Bader Ginsburg (541 U.S. at 534). Justice Ruth Bader Ginsburg filed a separate concurring opinion, joined by Justices David H. Souter and Stephen G. Breyer (541 U.S. at 535).
    \item \textbf{Dissents (4):} Chief Justice William H. Rehnquist filed a dissenting opinion, joined by Justices Anthony M. Kennedy and Clarence Thomas (541 U.S. at 538). Justice Antonin Scalia filed a separate dissenting opinion (541 U.S. at 554). Justice Clarence Thomas also filed a brief separate dissenting opinion (541 U.S. at 565) while joining Chief Justice Rehnquist's dissent in full.
\end{itemize}

\section{Holding}
\textbf{Judgment:} Affirmed. The Sixth Circuit's judgment permitting the Title~II money-damages suit to proceed against the State of Tennessee was affirmed (\href{https://www.law.cornell.edu/supremecourt/text/541/509}{541 U.S. at 534}).

\textbf{Rule Announced:} Under Section~5 of the Fourteenth Amendment and the ``congruence and proportionality'' framework of \textit{City of Boerne v. Flores}, 521 U.S. 507 (1997), Congress validly abrogated state sovereign immunity in Title~II of the ADA as it applies to the class of cases implicating the fundamental right of access to the courts (\href{https://www.law.cornell.edu/supremecourt/text/541/509}{541 U.S. at 533--534}). When evaluating a broad Section~5 statute that enforces multiple constitutional guarantees across diverse state programs, courts need not assess the statute ``as an undifferentiated whole,'' but may evaluate its congruence and proportionality as applied to the specific category of constitutional rights implicated by the case (\href{https://www.law.cornell.edu/supremecourt/text/541/509}{541 U.S. at 530--531}).

\textbf{Holding vs. Dicta:} The narrow holding is strictly limited to Title~II's validity under Section~5 ``as it applies to the class of cases implicating the fundamental right of access to the courts'' (\href{https://www.law.cornell.edu/supremecourt/text/541/509}{541 U.S. at 533--534}). The Court expressly declined to decide whether Title~II validly abrogates state sovereign immunity in other contexts---such as access to state-owned hockey rinks, voting booths, public education, or claims implicating only rational-basis equal protection review under \textit{City of Cleburne v. Cleburne Living Center, Inc.}, 473 U.S. 432 (1985) (\href{https://www.law.cornell.edu/supremecourt/text/541/509}{541 U.S. at 530--531, 532 n.20}). Accordingly, the majority's extensive survey of historical disability discrimination in voting, marriage, institutionalization, prisons, and public education serves as contextual support for Congress's legislative backdrop rather than a binding holding that Title~II abrogates sovereign immunity in those non-judicial settings (\href{https://www.cato.org/supreme-court-review/2003-2004/tennessee-v-lane-how-illegitimate-power-negated-non-existent-immunity}{Cato Supreme Court Review}).

\section{Analysis of Opinions}

\subsection{Majority Opinion (Justice Stevens)}
Justice Stevens applied the two-part abrogation framework from \textit{Kimel v. Florida Board of Regents}, 528 U.S. 62 (2000): first, whether Congress unequivocally expressed its intent to abrogate Eleventh Amendment immunity; and second, whether it acted pursuant to a valid grant of constitutional authority (\href{https://www.law.cornell.edu/supremecourt/text/541/509}{541 U.S. at 517}). The first prong was undisputedly met by 42 U.S.C. \S~12202 (\href{https://www.law.cornell.edu/supremecourt/text/541/509}{541 U.S. at 518}).

Under the second prong, Congress may abrogate state sovereign immunity when validly exercising its Fourteenth Amendment Section~5 enforcement power, \textit{Fitzpatrick v. Bitzer}, 427 U.S. 445 (1976), including by enacting prophylactic legislation that prohibits some facially constitutional conduct to deter or remedy constitutional violations, so long as the legislation exhibits ``congruence and proportionality'' between the injury and the remedial means under \textit{City of Boerne v. Flores}, 521 U.S. 507, 520 (1997) (\href{https://www.law.cornell.edu/supremecourt/text/541/509}{541 U.S. at 518--520}).

Distinguishing \textit{Garrett}, Justice Stevens emphasized that whereas Title~I enforced only the Equal Protection Clause's rational-basis ban on irrational disability discrimination in employment, Title~II also enforces a cluster of fundamental rights subject to heightened judicial scrutiny---specifically, the right of access to the courts protected by the Fourteenth Amendment Due Process Clause, the Sixth Amendment Confrontation and Jury Trial Clauses, and the First Amendment right of public access to criminal proceedings (\href{https://www.law.cornell.edu/supremecourt/text/541/509}{541 U.S. at 522--523}). Because infringements on fundamental rights trigger searching scrutiny, as in \textit{Nevada Department of Human Resources v. Hibbs}, 538 U.S. 721 (2003), it was easier for Congress to establish a pattern of constitutional violations (\href{https://www.law.cornell.edu/supremecourt/text/541/509}{541 U.S. at 528--529}).

Reviewing the historical record, Justice Stevens pointed to judicial decisions documenting unconstitutional treatment of disabled persons in state programs and judicial administration, a 1983 U.S. Commission on Civil Rights report finding that 76\% of public services in state-owned buildings were inaccessible, congressional hearing testimony on inaccessible courthouses, and extensive ADA Task Force submissions (\href{https://www.law.cornell.edu/supremecourt/text/541/509}{541 U.S. at 524--529}). He rejected Tennessee's argument that Title~II had to be evaluated facially across all public services, adopting instead a category-specific, as-applied inquiry focused on access to judicial services (\href{https://www.law.cornell.edu/supremecourt/text/541/509}{541 U.S. at 530--531}). Finally, Title~II's remedy was congruent and proportional because it requires only ``reasonable modifications'' that do not impose an ``undue financial or administrative burden'' or ``fundamental alteration,'' and allows older facilities to comply through lower-cost measures such as relocating proceedings or assigning aides, 28 C.F.R. \S~35.150 (\href{https://www.law.cornell.edu/supremecourt/text/541/509}{541 U.S. at 531--533}).

\subsection{Concurring Opinion (Justice Souter)}
Justice Souter, joined by Justice Ginsburg, joined the Court's opinion while preserving his broader reservations about the Court's Eleventh Amendment and Section~5 jurisprudence expressed in \textit{Hibbs} (\href{https://www.law.cornell.edu/supremecourt/text/541/509}{541 U.S. at 534}). Responding to Chief Justice Rehnquist's dissent, Justice Souter argued that a broader historical inquiry would further reinforce Congress's Section~5 authority because the judiciary itself had historically endorsed eugenicist laws imposing sterilization, institutionalization, and exclusions from schools, marriage, and voting on persons with disabilities---citing \textit{Buck v. Bell}, 274 U.S. 200 (1927), and \textit{State ex rel. Beattie v. Board of Education of Antigo}, 169 Wis. 231 (1919) (\href{https://www.law.cornell.edu/supremecourt/text/541/509}{541 U.S. at 534--535}).

\subsection{Concurring Opinion (Justice Ginsburg)}
Justice Ginsburg, joined by Justices Souter and Breyer, wrote separately to underscore the ADA's core equal-citizenship purpose (\href{https://www.law.cornell.edu/supremecourt/text/541/509}{541 U.S. at 535--536}). Citing \textit{Olmstead v. L.C.}, 527 U.S. 581 (1999), and disability-rights scholarship by Samuel Bagenstos, she explained that securing equal citizenship for persons with disabilities requires ``not blindfolded equality, but responsiveness to difference; not indifference, but accommodation'' (\href{https://www.law.cornell.edu/supremecourt/text/541/509}{541 U.S. at 536--537}). She also rejected Justice Scalia's suggestion that Congress must separately indict each individual State as a constitutional violator before enacting national Section~5 legislation, arguing that such an adversarial requirement is incompatible with a harmonious federal system (\href{https://www.law.cornell.edu/supremecourt/text/541/509}{541 U.S. at 537--538}).

\subsection{Dissenting Opinion (Chief Justice Rehnquist)}
Chief Justice Rehnquist, joined by Justices Kennedy and Thomas, argued that the majority's decision was irreconcilable with \textit{Garrett} and misapplied the three-step \textit{Boerne} framework (\href{https://www.law.cornell.edu/supremecourt/text/541/509}{541 U.S. at 538--540}). First, he contended that the legislative record was virtually devoid of evidence that States---as distinct from nonstate local governments, which lack Eleventh Amendment immunity---had engaged in a pattern of actual due process violations denying disabled persons access to the courts (\href{https://www.law.cornell.edu/supremecourt/text/541/509}{541 U.S. at 541--548}). Physical barriers alone do not violate the Due Process Clause when a State offers assistance or alternative venues, and neither the ADA's statutory findings nor its committee reports specifically identified courthouse access as a national constitutional crisis (\href{https://www.law.cornell.edu/supremecourt/text/541/509}{541 U.S. at 544--547}).

Second, Chief Justice Rehnquist sharply criticized the majority's ``as-applied'' methodology, arguing that it ``posits a hypothetical statute, never enacted by Congress, that applies only to courthouses'' and rigs the \textit{Boerne} test by artificially narrowing an undifferentiated statute that covers everything from courthouses to state-owned amusement parks and sports stadiums (\href{https://www.law.cornell.edu/supremecourt/text/541/509}{541 U.S. at 550--552}). Even confined to courthouses, he argued, Title~II imposes sweeping affirmative accommodation duties untethered to actual constitutional violations (\href{https://www.law.cornell.edu/supremecourt/text/541/509}{541 U.S. at 552--554}).

\subsection{Dissenting Opinion (Justice Scalia)}
Justice Scalia repudiated the \textit{Boerne} ``congruence and proportionality'' test altogether, calling it a ``flabby'' standard that invites ``judicial arbitrariness and policy-driven decisionmaking'' and forces the Court to act as ``Congress's taskmaster'' checking legislative homework (\href{https://www.law.cornell.edu/supremecourt/text/541/509}{541 U.S. at 557--558}). Turning to the text and 1860s dictionary meaning of ``enforce'' in Section~5, he argued that Congress may create causes of action and enforcement mechanisms for actual Fourteenth Amendment violations, but may not enact ``prophylactic'' legislation prohibiting primary conduct that the Constitution itself permits---\textit{i.e.}, ``reinforcement rather than enforcement'' (\href{https://www.law.cornell.edu/supremecourt/text/541/509}{541 U.S. at 558--560}). Out of deference to \textit{stare decisis} and the historical core of the Civil War Amendments, Justice Scalia stated he would retain permissive \textit{McCulloch v. Maryland} review only for Section~5 statutes targeting racial discrimination by the States (and only against States with an identified history of violations), while striking down all non-racial prophylactic legislation as \textit{ultra vires} (\href{https://www.law.cornell.edu/supremecourt/text/541/509}{541 U.S. at 560--565}).

\subsection{Dissenting Opinion (Justice Thomas)}
Justice Thomas joined Chief Justice Rehnquist's dissent in full and wrote a single-paragraph separate dissent solely to disavow any reliance on \textit{Nevada Department of Human Resources v. Hibbs}, 538 U.S. 721 (2003), which he continued to believe was wrongly decided (\href{https://www.law.cornell.edu/supremecourt/text/541/509}{541 U.S. at 565}).

\section{Examples: Future Applications}
\begin{enumerate}
    \item \textbf{Same-Side (Valid Abrogation):} A deaf civil litigant sues a State for money damages under Title~II after a state family court refuses to provide a sign-language interpreter for a parental-rights termination hearing. Because parental-rights termination implicates both fundamental due process access to the courts (\textit{M.L.B. v. S.L.J.}, 519 U.S. 102 (1996)) and the exact class of judicial services governed by \textit{Lane}, Congress validly abrogated sovereign immunity (\href{https://www.law.cornell.edu/supremecourt/text/541/509}{541 U.S. at 523, 533--534}).

    \item \textbf{Same-Side (Valid Abrogation):} A paraplegic state prison inmate sues the State for money damages under Title~II alleging that prison officials confined him in an inaccessible cell where he could not turn his wheelchair or access toilet facilities. As the Supreme Court later confirmed unanimously in \textit{United States v. Georgia}, 546 U.S. 151 (2006), insofar as Title~II creates a private damages remedy for conduct that \textit{actually} violates the Fourteenth Amendment (such as incorporated Eighth Amendment prohibitions against cruel and unusual punishment), Section~5 unquestionably authorizes abrogation (\href{https://supreme.justia.com/cases/federal/us/546/151/}{546 U.S. at 158--159}).

    \item \textbf{Opposite-Side (No Valid Abrogation):} A wheelchair user sues a State for money damages under Title~II because a state-owned public hockey rink or municipal-style state amusement park lacks accessible upper-deck seating. Unlike access to the courts, recreational spectator seating implicates no fundamental constitutional right and triggers only rational-basis scrutiny under \textit{Cleburne}; without a heightened constitutional right or a specific congressional record of unconstitutional state conduct in state recreation, \textit{Garrett} controls and sovereign immunity bars money damages (\href{https://www.law.cornell.edu/supremecourt/text/541/509}{541 U.S. at 530--531}; \href{https://www.cato.org/supreme-court-review/2003-2004/tennessee-v-lane-how-illegitimate-power-negated-non-existent-immunity}{Cato Supreme Court Review}).

    \item \textbf{Opposite-Side (No Valid Abrogation):} A blind applicant sues a state administrative agency for money damages under Title~II after being denied a state commercial fishing or barbering license without a requested accommodation, where the applicant alleges no procedural due process deprivation and the denial rests on rational-basis regulatory criteria. Because occupational licensing does not implicate a fundamental right and is reviewed under rational-basis scrutiny, lower courts applying \textit{Lane} and \textit{Garrett} generally hold that prophylactic Title~II damages remedies exceed Section~5 (\href{https://www.tascnow.com/wp-content/uploads/2019/03/TASC_0412_ada_lane.pdf}{TASC/NAPAS Post-Lane Survey}).

    \item \textbf{Fence-Sitter (Close Call):} A student with a learning disability sues a state university for money damages under Title~II after the university refuses to provide extended examination time, where the refusal does not independently violate the Equal Protection Clause. Although public higher education is not a ``fundamental right'' under \textit{San Antonio Independent School District v. Rodriguez}, 411 U.S. 1 (1973), \textit{Lane} highlighted a long history of disability discrimination in public education (\href{https://www.law.cornell.edu/supremecourt/text/541/509}{541 U.S. at 525 n.12}), leaving courts sharply divided over whether Congress's historical record in education suffices under \textit{Boerne} without a fundamental-right trigger (\href{https://scholarship.law.stjohns.edu/cgi/viewcontent.cgi?article=1203&context=lawreview}{St. John's Law Review}).
\end{enumerate}

\section{Critique}
\textbf{Originalist and Textualist Perspectives:} Originalist and textualist critics fault \textit{Lane} from two distinct angles. Echoing Justice Scalia's dissent, strict textualists argue that Section~5's grant of power ``to enforce'' the Fourteenth Amendment does not empower Congress to prohibit constitutionally permissible state conduct via prophylactic mandates, and that the \textit{Boerne} ``congruence and proportionality'' test is an atextual, indeterminate balancing standard that invites judicial policy-making (\href{https://www.law.cornell.edu/supremecourt/text/541/509}{541 U.S. at 557--559}; \href{https://uclawreview.org/2025/05/05/keeping-sovereign-immunity-kind-of-simple-why-courts-should-apply-the-united-states-v-georgia-test-strictly/}{University of Cincinnati Law Review}). 

\textbf{Progressive and Doctrinal Perspectives:} Progressive and disability-rights scholars generally welcomed \textit{Lane}'s outcome and Justice Ginsburg's recognition that equal citizenship requires affirmative accommodation, yet critiqued the majority's narrow, piecemeal ``as-applied'' holding (\href{http://openscholarship.wustl.edu/cgi/viewcontent.cgi?article=1332&context=law_journal_law_policy}{Washington University Journal of Law \& Policy}). By tethering Title~II's Section~5 validity to the fundamental right of court access rather than overturning \textit{Garrett} or treating disability as a suspect classification, \textit{Lane} invited lower courts to deny abrogation in vital non-judicial domains (such as public services, healthcare, and higher education) and forced disabled plaintiffs into protracted, category-by-category litigation over sovereign immunity (\href{http://openscholarship.wustl.edu/cgi/viewcontent.cgi?article=1332&context=law_journal_law_policy}{Washington University Journal of Law \& Policy}; \href{https://scholarship.law.stjohns.edu/cgi/viewcontent.cgi?article=1203&context=lawreview}{St. John's Law Review}).

\textbf{Competing Constitutional Values:} At bottom, \textit{Lane} reflects a deep structural clash among three constitutional commitments: (1)~federalism and state sovereign dignity (protecting state treasuries from federally mandated damages liability), (2)~separation of powers and judicial supremacy (preventing Congress from using Section~5 to redefine the substantive scope of Fourteenth Amendment rights), and (3)~the Reconstruction Amendments' transformative promise of national civil-rights enforcement to dismantle entrenched state-level exclusion of marginalized groups (\href{https://www.law.cornell.edu/supremecourt/text/541/509}{541 U.S. at 518--520, 536--538}).

\section{Key Quotations}
\begin{quote}
``For these reasons, we conclude that Title II, as it applies to the class of cases implicating the fundamental right of access to the courts, constitutes a valid exercise of Congress' \S~5 authority to enforce the guarantees of the Fourteenth Amendment.''
\end{quote}
\noindent --- Justice John Paul Stevens, Opinion of the Court, 541 U.S. at 533--534 (\href{https://www.law.cornell.edu/supremecourt/text/541/509}{Cornell LII}; \href{https://www.courtlistener.com/opinion/134743/tennessee-v-lane/}{CourtListener}).

\begin{quote}
``Including individuals with disabilities among people who count in composing `We the People,' Congress understood in shaping the ADA, would sometimes require not blindfolded equality, but responsiveness to difference; not indifference, but accommodation.''
\end{quote}
\noindent --- Justice Ruth Bader Ginsburg, Concurring Opinion, 541 U.S. at 536 (\href{https://www.law.cornell.edu/supremecourt/text/541/509}{Cornell LII}).

\begin{quote}
``In conducting its as-applied analysis, however, the majority posits a hypothetical statute, never enacted by Congress, that applies only to courthouses. The effect is to rig the congruence-and-proportionality test by artificially constricting the scope of the statute to closely mirror a recognized constitutional right.''
\end{quote}
\noindent --- Chief Justice William H. Rehnquist, Dissenting Opinion, 541 U.S. at 551 (\href{https://www.law.cornell.edu/supremecourt/text/541/509}{Cornell LII}).

\begin{quote}
``The `congruence and proportionality' standard, like all such flabby tests, is a standing invitation to judicial arbitrariness and policy-driven decisionmaking. Worse still, it casts this Court in the role of Congress's taskmaster.''
\end{quote}
\noindent --- Justice Antonin Scalia, Dissenting Opinion, 541 U.S. at 557--558 (\href{https://www.law.cornell.edu/supremecourt/text/541/509}{Cornell LII}).

\begin{quote}
``Nothing in \S~5 allows Congress to go beyond the provisions of the Fourteenth Amendment to proscribe, prevent, or `remedy' conduct that does not itself violate any provision of the Fourteenth Amendment. So-called `prophylactic legislation' is reinforcement rather than enforcement.''
\end{quote}
\noindent --- Justice Antonin Scalia, Dissenting Opinion, 541 U.S. at 559 (\href{https://www.law.cornell.edu/supremecourt/text/541/509}{Cornell LII}).


\end{document}
